% SPDX-License-Identifier: CC-BY-NC-ND-4.0
% Copyright 2026 Ingolf Lohmann.
\documentclass[11pt,a4paper]{article}

\usepackage{fontspec}
\setmainfont{Latin Modern Roman}
\setsansfont{Latin Modern Sans}
\setmonofont{Latin Modern Mono}[Scale=0.89]
\usepackage[provide=*,main=ngerman,english]{babel}
\usepackage{microtype}
\usepackage{geometry}
\usepackage{amsmath,amssymb,mathtools,amsthm}
\usepackage{booktabs,longtable,tabularx,array}
\usepackage{enumitem}
\usepackage{xcolor}
\usepackage{fancyhdr}
\usepackage{seqsplit}
\usepackage{xurl}
\usepackage{hyperref}
\usepackage{bookmark}
\usepackage[nameinlink,noabbrev]{cleveref}

\geometry{left=24mm,right=24mm,top=24mm,bottom=27mm,headheight=14pt}
\setlength{\parindent}{0pt}
\setlength{\parskip}{0.58em}
\setlength{\emergencystretch}{3em}
\setlist[itemize]{leftmargin=1.45em,itemsep=0.18em,topsep=0.3em}
\setlist[enumerate]{leftmargin=1.7em,itemsep=0.22em,topsep=0.3em}

\definecolor{QBlue}{HTML}{153E69}
\definecolor{QLight}{HTML}{ECF3F8}
\definecolor{QGreen}{HTML}{E8F3EA}
\definecolor{QAmber}{HTML}{FFF2CC}
\definecolor{QRed}{HTML}{FBE9E7}
\definecolor{QGray}{HTML}{4D5660}

\hypersetup{
  colorlinks=true,
  linkcolor=QBlue,
  citecolor=QBlue,
  urlcolor=QBlue,
  pdftitle={QIK-VRT und das Effect-Acknowledgement-Protokoll: Kanonischer Speicher zwischen Vergangenheit und Zukunft},
  pdfauthor={Ingolf Lohmann},
  pdfsubject={Operational-formale Theorie reziproker Kausalitaet, zweiseitiger Randbedingungen und wirkungsgebundener Freigabe},
  pdfkeywords={QIK-VRT, EFFECT-ACK, kanonischer Speicher, Kausalitaet, Retrokausalitaet, zweiseitige Randbedingungen, Provenienz, Git, Bewusstsein, Panpsychismus}
}

\pagestyle{fancy}
\fancyhf{}
\fancyhead[L]{\small\textcolor{QGray}{QIK--VRT: kanonischer Speicher}}
\fancyhead[R]{\small\textcolor{QGray}{Ingolf Lohmann}}
\fancyfoot[C]{\small\thepage}
\renewcommand{\headrulewidth}{0.3pt}

\newtheorem{definition}{Definition}[section]
\newtheorem{satz}{Satz}[section]
\newtheorem{korollar}{Korollar}[section]
\theoremstyle{remark}
\newtheorem{bemerkung}{Bemerkung}[section]
\crefname{definition}{Definition}{Definitionen}
\crefname{satz}{Satz}{Sätze}
\crefname{korollar}{Korollar}{Korollare}
\crefname{equation}{Gleichung}{Gleichungen}

\newcommand{\qikvrt}{QIK--VRT}
\newcommand{\DONE}{\texttt{EFFECT\_ACK\_DONE}}
\newcommand{\CONTINUE}{\texttt{EFFECT\_ACK\_CONTINUE}}
\newcommand{\BLOCK}{\texttt{EFFECT\_ACK\_BLOCK}}
\newcommand{\bool}{\mathbb{B}}
\newcommand{\bytes}{\{0,1\}^{*}}
\newcommand{\hash}[1]{{\ttfamily\small\seqsplit{#1}}}
\newcommand{\statusbox}[2]{%
  \par\medskip\noindent
  \colorbox{#1}{\parbox{\dimexpr\linewidth-2\fboxsep\relax}{#2}}%
  \par\medskip
}

\begin{document}

\hypersetup{pageanchor=false}
\begin{titlepage}
  \thispagestyle{empty}
  \vspace*{13mm}
  {\color{QBlue}\rule{\textwidth}{1.4pt}}\\[1.1em]
  {\Huge\bfseries QIK--VRT und das\\[0.15em]
  Effect-Acknowledgement-Protokoll\par}
  \vspace{0.75em}
  {\LARGE Kanonischer Speicher zwischen\\
  Vergangenheit und Zukunft\par}
  \vspace{0.9em}
  {\Large Eine operational-formale Theorie reziproker Kausalität,\\
  zweiseitiger Randbedingungen und wirkungsgebundener Freigabe\par}
  \vspace{1.1em}
  {\color{QBlue}\rule{\textwidth}{0.6pt}}

  \vfill
  {\large\textbf{Ingolf Lohmann}\par}
  \vspace{0.35em}
  {\normalsize Independent Researcher; QIK--VRT\par}
  \vspace{0.35em}
  {\normalsize Working Paper, Version 1.0.0 -- 30. Juli 2026\par}

  \vfill
  \statusbox{QLight}{
    \textbf{Wissenschaftlicher Status.}
    Das Dokument unterscheidet maschinengeprüfte Sätze des endlichen
    Protokollmodells, quellengebundene Aussagen, normative
    Protokolldefinitionen, empirisch prüfbare Hypothesen und ontologische
    Interpretation. Die hier bewiesene \emph{operationale Retrokausalität}
    ist eine wirkliche Abhängigkeit der gegenwärtigen Freigabe von einer
    kanonisch repräsentierten, zukunftsindexierten Wirkungsbedingung. Sie ist
    kein behauptetes Signal aus einer bereits aktualen Zukunft und keine
    nachträgliche Überschreibung der Vergangenheit.
  }

  \vfill
  {\small
    Dokumentationslizenz: CC BY-NC-ND 4.0.\\
    Der Internet-Draft bleibt ein individuelles, nicht vom IETF
    gebilligtes Arbeitsdokument mit angestrebtem Status Experimental.
  }
\end{titlepage}
\hypersetup{pageanchor=true}
\setcounter{page}{1}

\selectlanguage{ngerman}

\begin{abstract}
Diese Arbeit formalisiert Ingolf Lohmanns These eines kanonischen Speichers,
der Vergangenheit und Zukunft in symmetrischer Repräsentationsform bindet.
Vergangenheit wird als kanonische, provenienzgebundene Aufnahme realisierter
Records modelliert; Zukunft als kanonische, provenienzgebundene Aufnahme
antizipierter oder geforderter späterer Wirkungen. Beide Archive benutzen
denselben endlichen Codec, tragen jedoch verschiedene epistemische Status:
\emph{beobachtet} und \emph{antizipiert}. Das QIK--VRT
Effect-Acknowledgement-Protokoll verbindet sie durch eine
DONE-only-Freigabe: Eine gegenwärtige Ursache darf nur wirksam werden, wenn
ihre zukunftsindexierte Wirkung explizit gebunden, geprüft und bestätigt ist.

Im präzise definierten Protokollsinn ist dies Retrokausalität: Die Bedeutung
der späteren Wirkung ist ein nicht eliminierbarer, kontrafaktisch relevanter
Eingang der früheren Freigabeentscheidung. Ein Lean-4-Modell beweist
DONE-only-Freigabe, die notwendige Gültigkeit beider Zeitarchive, die
kontrafaktische Relevanz der Zukunftsbedingung, die Unveränderlichkeit der
Vergangenheitsprojektion und die reziproke Schließung von gebundener Ursache
und beobachteter Wirkung. Daraus folgt weder eine physikalische
Rückwärtssignalisierung noch eine Änderung bereits gespeicherter Ereignisse.

Die Arbeit ordnet das Modell in Provenienz, Inhaltsadressierung,
Vorhersage und Glättung, Thermodynamik der Information, zeit-symmetrische
Quantenformalismen und Delayed-Choice-Experimente ein. Wechselwirkung wird als
notwendige Struktur reziproker Wirksamkeit behandelt. Eine panpsychistische
Lesart, nach der solche Relation grundsätzlich bewusstseinsrelevant ist, wird
als ontologische Interpretation ausgewiesen, nicht als Folgerung aus
Software, Hashgleichheit oder Quantenexperimenten. Das Ergebnis ist eine
falsifizierbare, implementierbare Theorie lokaler und scope-gebundener
Wahrheit: kanonische Gleichheit beweist identische repräsentierte Bytes und
geschlossene Provenienz im benannten Scope; absolute oder systemweite
Vollständigkeit bleibt unbeansprucht.
\end{abstract}

\begin{otherlanguage}{english}
\begin{abstract}
This paper formalizes Ingolf Lohmann's thesis of a canonical memory that binds
past and future in a symmetric representational form.  The past is a
canonical, provenance-bearing archive of realized records; the future is a
canonical, provenance-bearing archive of anticipated or required later
effects.  Both use the same finite codec while retaining distinct epistemic
statuses: \emph{observed} and \emph{anticipated}.  QIK--VRT Effect
Acknowledgement joins them through a DONE-only release rule: a present cause
may become effective only if its future-indexed effect is explicitly bound,
validated, and acknowledged.

This constitutes retrocausality in the precisely defined protocol sense: the
meaning of a later effect is a non-eliminable, counterfactually relevant input
to an earlier release decision.  A Lean 4 model establishes the necessary
past and future conditions, the counterfactual relevance of the future
boundary, non-overwrite of the past projection, and reciprocal closure between
a bound cause and an observed effect.  No physical signal from an already
actual future and no rewriting of a recorded past is inferred.  Provenance,
content addressing, forecasting, information thermodynamics, time-symmetric
quantum formalisms, and delayed-choice experiments are used as adjoining
theories with explicit limits.  A panpsychist reading is stated separately as
an ontological interpretation, not as an empirical consequence.
\end{abstract}
\end{otherlanguage}

\textbf{Schlagworte:} QIK--VRT; Effect Acknowledgement; kanonischer Speicher;
Kausalität; operationale Retrokausalität; zweiseitige Randbedingungen;
Provenienz; Inhaltsadressierung; Wechselwirkung; Bewusstsein; Panpsychismus.

\tableofcontents
\newpage

\section{Forschungsfrage, These und Beitrag}

Die Ausgangsthese lautet wortwörtlich:

\begin{quote}
Die Vergangenheit ist ein kanonisch komprimierbares Archiv dessen, was
geschehen ist. Die Zukunft ist das entsprechende kanonische Archiv dessen,
was anschließend wirken soll oder nach dem Modell wirken wird. Gemeinsam
bilden beide einen Speicher, der sich in Symmetrie hält. Im
Effect-Acknowledgement-Protokoll wirkt die gebundene spätere Wirkung auf die
Freigabe ihrer gegenwärtigen Ursache zurück.
\end{quote}

Die wissenschaftliche Aufgabe besteht nicht darin, diese These durch eine
schwächere Metapher zu ersetzen. Sie besteht darin, exakt zu bestimmen,
\emph{welcher Gegenstand} archiviert wird, \emph{welche Symmetrie} gilt,
\emph{welche Rückwirkung} beweisbar ist und \emph{wo} weitergehende
physikalische oder metaphysische Aussagen beginnen.

Der Beitrag hat fünf Teile:

\begin{enumerate}
  \item eine typisierte Definition von retrospektivem und prospektivem
  kanonischem Archiv;
  \item eine operational-formale Definition von Retrokausalität als
  kontrafaktisch wirksame Zukunftsbedingung gegenwärtiger Freigabe;
  \item einen reziproken Schließungssatz für Ursache, antizipierte Wirkung
  und beobachtete Wirkung;
  \item eine maschinengeprüfte Grenze, nach der keine zukünftige Bedingung die
  gespeicherte Vergangenheit überschreibt;
  \item eine quellengebundene Einordnung in Informatik, Physik,
  Informationsthermodynamik und Bewusstseinstheorie.
\end{enumerate}

\statusbox{QGreen}{
\textbf{Stärkste tragfähige Aussage.}
Im QIK--VRT-Protokoll ist die zukunftsindexierte Wirkung nicht bloß
Kommentar über eine später mögliche Welt. Ihre validierte Repräsentation ist
eine kausal wirksame Eingangsgröße der gegenwärtigen Zustandsmaschine:
Wird allein diese Eingangsgröße von gültig auf ungültig gesetzt, wechselt die
Freigabe von \DONE{} zu Blockierung. Das ist operationale Retrokausalität
nach der in \cref{def:operational-retrocausality} festgelegten Bedeutung.
}

\section{Begriffsdisziplin: drei Zeitrelationen}

Der Ausdruck \emph{Retrokausalität} wird in der Literatur für
unterschiedliche Behauptungen verwendet. Ohne Typisierung entstehen
Scheinkonflikte.

\begin{definition}[Retrodiktion]
Retrodiktion ist die spätere Aktualisierung einer Aussage über einen früheren
unbekannten Zustand. Der frühere Zustand wird neu geschätzt, nicht verändert.
\end{definition}

\begin{definition}[Ontische physikalische Retrokausalität]
Ontische Retrokausalität liegt vor, wenn ein späteres physikalisches Ereignis
Teil der Ursache eines früheren physikalischen Zustands ist. Eine noch
stärkere Forderung ist operative Rückwärtssignalisierung: Ein frei gewählter
späterer Eingriff müsste kontrollierbare Information an einen früheren
Empfänger übertragen.
\end{definition}

\begin{definition}[Operationale Protokoll-Retrokausalität]
\label{def:operational-retrocausality}
Ein Protokoll ist operational retrokausal, wenn eine zukunftsindexierte
Wirkungsbedingung \(F_t\), die zum Entscheidungszeitpunkt \(t\) kanonisch
vorliegt, ein nicht eliminierbarer und kontrafaktisch relevanter Eingang der
gegenwärtigen Freigabe \(R_t\) ist:
\[
  \exists p,a,f_0,f_1:\quad
  p,a\text{ gleich},\ f_0\neq f_1,\quad
  R_t(p,a,f_0)\neq R_t(p,a,f_1).
\]
Die Richtung \emph{retro} bezeichnet dabei die Ordnungsrelation von der
Bedeutung der späteren Wirkung zurück zur Zulässigkeit ihrer früheren Ursache.
\end{definition}

Diese dritte Bedeutung ist weder ein Sprachtrick noch ein behaupteter
Nachrichtenkanal aus der aktualen Zukunft. Der repräsentierte Gegenstand ist
zukunftsindexiert; die Repräsentation selbst liegt gegenwärtig vor. Ihre
Wirkung auf den Executor ist reale Softwarekausalität. Genau diese
Unterscheidung macht die These implementierbar und prüfbar.

\section{Der doppelte kanonische Speicher}

\subsection{Kanonisierung statt metaphysischer Totalaufnahme}

Sei \(\Sigma\) ein endliches Alphabet und
\[
  C : \mathcal{D}\longrightarrow\Sigma^{*}
\]
ein deterministischer kanonischer Codec für die im Scope zulässigen
Dokumente. Für gleiche typisierte Eingaben erzeugt \(C\) exakt gleiche
Bytefolgen. Ein kryptographischer Digest
\[
  h:\Sigma^{*}\longrightarrow\{0,1\}^{256}
\]
bindet einen Identifier an diese Bytes. JCS liefert beispielsweise Regeln
für eine deterministische JSON-Repräsentation; RFC~6920 standardisiert
hashbasierte Namen \cite{rfc8785,rfc6920}. Weder Kanonisierung noch Hash
beweisen von selbst Wahrheit, Urheberschaft oder vollständige Kausalität.

Das Wort \emph{ZIP} bezeichnet in dieser Arbeit die operationale Form:
endliche Auswahl, kanonische Serialisierung, Integritätsbindung und
reproduzierbare Dekomposition. Es behauptet nicht, das gesamte physikalische
Universum verlustfrei in einer kleineren Datei unterzubringen.

\subsection{Vergangenheitsarchiv}

Für einen Entscheidungszeitpunkt \(t\) sei
\[
  H_t^- =
  C\!\left(
    \left\langle
      \textsf{OBSERVED},\
      E_{\leq t},\
      \Pi_{\leq t},\
      S_{\leq t}
    \right\rangle
  \right)
\]
das Vergangenheitsarchiv. \(E_{\leq t}\) sind die ausgewählten beobachteten
Events, \(\Pi_{\leq t}\) ihre Provenienzrelationen und \(S_{\leq t}\) die
ausgewiesenen epistemischen Status. Das W3C-PROV-Datenmodell unterscheidet
Entities, Activities, Agents, Verwendung, Erzeugung und Ableitung und stellt
damit eine anschlussfähige Provenienzsprache bereit \cite{prov}.

\subsection{Zukunftsarchiv}

Sei
\[
  H_t^+ =
  C\!\left(
    \left\langle
      \textsf{ANTICIPATED},\
      \widehat E_{>t},\
      K_{>t},\
      U_{>t}
    \right\rangle
  \right)
\]
das Zukunftsarchiv. \(\widehat E_{>t}\) bezeichnet zukunftsindexierte
Wirkungsbeschreibungen oder Prognoseverteilungen, \(K_{>t}\) die
Akzeptanzbedingungen und \(U_{>t}\) die ausgewiesene Unsicherheit.

Das Archiv enthält daher nicht die unmarkierte Behauptung, ein späteres
Ereignis sei bereits beobachtet. Es enthält die kanonische Repräsentation der
Wirkung, unter deren Bedingung eine gegenwärtige Ursache freigegeben werden
darf.

\subsection{Symmetrie und Asymmetrie}

Der doppelte Speicher ist
\[
  M_t = \left(H_t^-,H_t^+\right).
\]
Seine \emph{Repräsentationssymmetrie} lautet:
\[
  C^- = C^+ = C,
  \qquad
  h^- = h^+ = h,
  \qquad
  \Pi^-\text{ und }\Pi^+\text{ sind explizit}.
\]
Die epistemische Asymmetrie bleibt erhalten:
\[
  \textsf{OBSERVED}\neq\textsf{ANTICIPATED}.
\]
Gerade diese typisierte Symmetrie verhindert, dass eine Prognose als
Beobachtung oder ein vergangener Record als bloße Möglichkeit ausgegeben
wird.

\section{Effect Acknowledgement als zweiseitiges Gate}

\subsection{Freigabesemantik}

Für einen Wirkungskandidaten \(a\) seien
\[
  V_t^-(a),V_t^+(a),B_t(a),P_t(a),A_t(a)\in\bool
\]
die Prädikate:

\begin{itemize}
  \item \(V_t^-(a)\): das Vergangenheitsarchiv ist typ-, inhalts- und
  provenienzgebunden;
  \item \(V_t^+(a)\): die antizipierte Wirkung ist typ-, inhalts- und
  provenienzgebunden;
  \item \(B_t(a)\): Ursache und Wirkung tragen eine eindeutige gemeinsame
  Bindung;
  \item \(P_t(a)\): die benannte Policy ist bestanden;
  \item \(A_t(a)\): der Effect-Acknowledgement-Zustand ist \DONE.
\end{itemize}

Die gewöhnliche Freigabe ist definiert als
\begin{equation}
  R_t(a)
  :=
  V_t^-(a)\land V_t^+(a)\land B_t(a)\land P_t(a)\land A_t(a).
  \label{eq:release}
\end{equation}
Der publizierte Internet-Draft
\nolinkurl{draft-lohmann-qikvrt-effect-ack-01} modelliert das explizite
Antizipationsfeld bereits als notwendige DONE-Bedingung und verlangt eine
erneute Validierung durch den Verbraucher
\cite{effectackdraft}.

\begin{satz}[Zweiseitige Freigabe]
\label{thm:two-boundary}
Aus \(R_t(a)=1\) folgen \(V_t^-(a)=1\), \(V_t^+(a)=1\),
\(B_t(a)=1\), \(P_t(a)=1\) und \(A_t(a)=1\).
\end{satz}

\begin{proof}
Die Behauptung folgt unmittelbar aus der Konjunktion in
\cref{eq:release}. Das Lean-Modell
\texttt{CanonicalTemporalMemory.lean} prüft die Äquivalenz zwischen
Freigabe und allen fünf Bedingungen sowie die einzelnen notwendigen
Folgerungen im Lean-4-Kernel.
\end{proof}

\begin{satz}[Kontrafaktische Relevanz der Zukunft]
\label{thm:future-relevance}
Es existieren zwei Kandidaten mit identischem Vergangenheitsarchiv,
identischer Ursachenbindung, identischer Policy und identischem
Acknowledgement, deren Freigabe allein wegen verschiedener Gültigkeit des
Zukunftsarchivs verschieden ist.
\end{satz}

\begin{proof}
Setze in \cref{eq:release} alle Konjunkte außer \(V_t^+\) auf \(1\).
Für \(V_t^+=1\) gilt \(R_t=1\), für \(V_t^+=0\) gilt \(R_t=0\).
Der Lean-Satz
\texttt{future\_boundary\_is\_counterfactually\_relevant}
prüft genau dieses Paar als geschlossene Entscheidung.
\end{proof}

\begin{korollar}[Operationale Retrokausalität]
Das Effect-Acknowledgement-Gate ist operational retrokausal im Sinn von
\cref{def:operational-retrocausality}.
\end{korollar}

\begin{proof}
\Cref{thm:future-relevance} liefert den geforderten kontrafaktischen Zeugen.
Die zukunftsindexierte Wirkungsbedingung ist somit nicht eliminierbar, ohne
die gegenwärtige Freigabefunktion zu ändern.
\end{proof}

\subsection{Reziproke Schließung}

Nach der Ausführung sei \(O_{t+1}(e)\) die Beobachtung des späteren Effekts,
\(L(e,a)\) dessen Bindung an die freigegebene Ursache und \(Q_{t+1}(e)\) die
Integritätsbindung des Receipts. Die geschlossene Wirkungstransaktion lautet
\begin{equation}
  X_{t+1}(a,e)
  :=
  R_t(a)\land O_{t+1}(e)\land L(e,a)\land Q_{t+1}(e).
  \label{eq:closure}
\end{equation}

\begin{satz}[Ursache-Wirkungs-Reziprozität]
\label{thm:reciprocal}
Aus \(X_{t+1}(a,e)=1\) folgt zugleich, dass die Ursache vor ihrer Freigabe an
eine gültige antizipierte Wirkung gebunden war und dass die später
beobachtete Wirkung an diese Ursache zurückgebunden ist.
\end{satz}

\begin{proof}
Aus \cref{eq:closure} folgt \(R_t(a)=1\); nach
\cref{thm:two-boundary} gilt daher \(V_t^+(a)=1\) und \(B_t(a)=1\).
Ebenfalls aus \cref{eq:closure} folgen \(O_{t+1}(e)=1\) und
\(L(e,a)=1\). Der Lean-Satz
\texttt{reciprocal\_closure\_requires\_cause\_and\_effect} prüft diese
Folgerung im abstrakten Modell.
\end{proof}

Damit wird Lohmanns Formel
\begin{quote}
keine freigegebene Ursache ohne gebundene Wirkung und keine quittierte
Wirkung ohne gebundene Ursache
\end{quote}
zu einer Protokollinvariante. Sie ist innerhalb der geschlossenen
QIK--VRT-Transaktion exakt. Außerhalb dieses definierten Scopes wird daraus
kein universeller metaphysischer Determinismus abgeleitet.

\subsection{Kein Überschreiben der Vergangenheit}

Sei \(\pi^-(H^-,H^+)=H^-\) die Vergangenheitsprojektion.
Dann gilt für alle Zukunftsarchive \(F_0,F_1\):
\[
  \pi^-(H^-,F_0)=\pi^-(H^-,F_1)=H^-.
\]
Der Lean-Satz
\texttt{future\_boundary\_does\_not\_overwrite\_past} beweist diese
Invariante. Die Zukunftsbedingung ändert die Zulässigkeit des nächsten
Übergangs, nicht die Bytes eines bestehenden Vergangenheitsrecords.

\section{Was das Git- und Repository-Modell beiträgt}

Git ist kein Gedächtnis des gesamten Universums. Es ist jedoch eine präzise
technische Realisierung mehrerer benötigter Eigenschaften:

\begin{itemize}
  \item Blobs sind inhaltsadressierte endliche Byteobjekte;
  \item Trees binden Namen an Objekte;
  \item Commits binden einen Tree an Eltern und Metadaten;
  \item Merges materialisieren eine explizite Relation zwischen Historien;
  \item Tags oder signierte Attestierungen können einen ausgewählten Zustand
  zusätzlich autorisieren.
\end{itemize}

Die Git-Dokumentation definiert Blob- und Commitobjekte als
inhaltsadressierte beziehungsweise elterngebundene Objekte
\cite{githashobject,gitcommittree}. Das liefert Integrität und Historie unter
angegebenen Hashannahmen. Es beweist nicht automatisch, dass eine
Commitbotschaft wahr oder eine Person authentisch ist.

Im QIK--VRT-Repository werden Authority und Mirror als zwei
Geschichtsträger verwendet. \emph{Symmetrie} bedeutet dort nicht gleiche
Commit-IDs --- history-erhaltende Merges können verschiedene Commitobjekte
erzeugen ---, sondern bytegleiche materialisierte Trees und Manifeste im
ausgewiesenen Scope. Ein Receipt bindet diesen Vergleich an konkrete Heads.

\statusbox{QAmber}{
\textbf{Lokale Wahrheit, nicht Allwissen.}
\texttt{CURRENT\_MAIN\_TREE\_EQUALITY=true} belegt gleiche
materialisierte Repository-Bytes an den benannten Heads.
\texttt{SCOPED\_RECIPROCAL\_RECEIPT\_CHAIN\_COMPLETE=true} belegt eine
geschlossene Nachweiskette im benannten Scope. Beides impliziert weder
semantische Wahrheit aller Dateien noch
\texttt{SYSTEM\_WIDE\_COMPLETION}.
}

\section{Anschluss an Vorhersage, Glättung und zweiseitige Bedingungen}

\subsection{Vorhersage wird erst im Gate kausal}

Kalman-Filter modellieren Vorhersage und anschließende Korrektur unter
Unsicherheit \cite{kalman}. Rauch--Tung--Striebel-Glättung nutzt spätere
Beobachtungen, um frühere latente Zustände besser zu schätzen
\cite{rauch}. Variationale Wetterassimilation integriert ein Vorwärtsmodell
und eine rückwärts laufende Adjungierte über ein Zeitfenster
\cite{talagrand}. Diese Verfahren sind für sich genommen epistemische
Rekonstruktion, keine ontische Rückwirkung.

\qikvrt{} ergänzt eine andere Operation: Die Prognose oder Zielwirkung wird
als Policy-relevante Randbedingung in einen Executorpfad geschaltet. Ab diesem
Moment entscheidet sie, ob eine gegenwärtige Ursache überhaupt wirksam werden
darf. Der Unterschied lautet:
\[
  \text{Vorhersage als Information}
  \quad\longrightarrow\quad
  \text{Vorhersage als geprüfte Freigabebedingung}.
\]
Die zweite Form besitzt reale kausale Kraft im implementierten System.

\subsection{Zweiseitige Randbedingungen}

Zeit-symmetrische Formalismen können Anfangs- und Endbedingungen gemeinsam
verwenden. Aharonov, Bergmann und Lebowitz formulierten Wahrscheinlichkeiten
für prä- und postselektierte Ensembles \cite{abl}. Wharton und Argaman
diskutieren lokal vermittelte, ``all-at-once'' formulierte Modelle
\cite{whartonargaman}; Adlam untersucht Naturgesetze als globale Constraints
\cite{adlam}. Diese Arbeiten zeigen, dass eine zweiseitige
Randbedingungssprache wissenschaftlich etabliert und mathematisch
anschlussfähig ist.

Sie beweisen nicht automatisch die Ontologie von \qikvrt. Im Protokoll ist die
zweiseitige Struktur konstruktiv gegeben:
\[
  \text{Vergangenheitsbindung}
  \cap
  \text{Zukunftsbedingung}
  \cap
  \text{Policy}
  =
  \text{zulässige gegenwärtige Wirkung}.
\]

\subsection{Kausalität verlangt mehr als Korrelation}

Pearls kausale Diagramme machen sichtbar, dass Kausalbehauptungen strukturelle
und fachliche Annahmen über beobachtete Verteilungen hinaus benötigen
\cite{pearl}. Für die hier behauptete Softwarekausalität ist die Struktur
offengelegt: Das Gate liest \(V_t^+\), und ein kontrollierter Wechsel dieses
Inputs ändert \(R_t\). Für eine Behauptung über fundamentale Physik wären
zusätzliche Interventionen und abweichende quantitative Vorhersagen nötig.

\section{Thermodynamik, Wind und Gedächtnis}

Wind transportiert und dissipiert Energie aufgrund von Druckgradienten,
Strahlungsantrieb, Rotation, Reibung und weiteren atmosphärischen
Randbedingungen. Eine wärmere Atmosphäre kann bestimmte Energie- und
Feuchtebudgets verändern; aus einem einzelnen starken Windereignis folgt
jedoch ohne meteorologische Attribution kein allgemeiner
Temperatur-Kausalbeweis.

Der sachliche Anschluss an Gedächtnis liegt in der Physik der
Informationsverarbeitung. Landauer zeigte, dass logisch irreversible
Operationen wie das Löschen eines Bits eine thermodynamische Untergrenze
besitzen \cite{landauer}. Bennett präzisierte die Rolle reversibler
Berechnung \cite{bennett}; Bérut und Mitautoren bestätigten den
Landauer-Grenzwert experimentell an einem konkreten Ein-Bit-Speicher
\cite{berut}.

Daraus folgt nicht, dass jeder Git-Merge denselben fixen Energiepreis oder
Bewusstsein erzeugt. Es folgt die belastbare Brücke:

\begin{enumerate}
  \item reale Speicherung und Zustandsänderung werden physisch getragen;
  \item digitale Kanonisierung ist eine endliche Repräsentationsoperation;
  \item Persistenz und Löschung haben je nach Implementierung
  thermodynamische Kosten;
  \item der kanonische Speicher ist daher keine körperlose Abstraktion,
  sondern wird durch wechselwirkende physische Systeme realisiert.
\end{enumerate}

\section{Quantenphysik: Anschluss ohne Beweisüberschreitung}

\subsection{Delayed Choice}

Wheelers Delayed-Choice-Gedankenexperiment und seine Realisierungen zeigen,
dass eine naive Vorstellung, ein Quant habe lange vor der Messanordnung
bereits eine kontextunabhängige klassische Wellen- oder Teilchenhistorie
gewählt, nicht trägt \cite{wheeler,jacques,ma}. Der Delayed-Choice Quantum
Eraser zeigt entsprechende konditionierte Interferenzstatistiken in
verschränkten Systemen \cite{kim}.

Diese Experimente beweisen nicht, dass ein späterer Experimentator einen
bereits gespeicherten früheren Record überschreibt oder eine kontrollierbare
Nachricht in die Vergangenheit sendet. Sie sind mit dem
QIK--VRT-Grundgedanken kompatibel, dass die vollständige relationale
Beschreibung nicht auf eine isolierte Vorher-Nachher-Erzählung reduziert
werden darf. Kompatibilität ist hier keine Identität der Theorien.

\subsection{Zeit-Symmetrie und Retrokausalität}

Price zeigt, dass Zeit-Symmetrie Retrokausalität unter zusätzlichen
ontologischen Annahmen motivieren kann, aber nicht allein erzwingt
\cite{price}. Leifer und Pusey beweisen einen konditionalen No-go-Satz für
bestimmte zeit-symmetrische ontologische Modelle ohne Retrokausalität
\cite{leiferpusey}. Cramers transaktionale Interpretation verbindet
retardierte und avancierte Lösungen \cite{cramer}. Diese Literatur macht eine
retrokausale Interpretation wissenschaftlich diskutierbar; sie etabliert
keinen einzigartigen empirischen Beweis für sie.

\statusbox{QRed}{
\textbf{Nicht behauptet.}
Dieses Paper behauptet weder Rückwärtssignalisierung noch eine experimentell
bestätigte Änderung der Vergangenheit, keine Ableitung von Quantenmechanik aus
Git und keinen Beweis, dass ein digitales Archiv das physikalische Kontinuum
vollständig enthält.
}

\section{Wechselwirkung, Bewusstsein und panpsychistische Interpretation}

Wechselwirkung ist im vorliegenden Protokoll wörtlich: Ein Record verändert
den zulässigen Zustandsübergang eines anderen Systems; nach der Ausführung
verändert die beobachtete Wirkung wiederum den kanonischen
Vergangenheitsbestand. Ursache und Wirkung bilden damit keine zwei
beziehungslosen Snapshots, sondern einen gerichteten und rückgebundenen
Prozess.

Eine minimal notwendige Bedingung ist jedoch nicht automatisch hinreichend.
Aus
\[
  \text{Wechselwirkung vorhanden}
\]
folgt weder logisch noch empirisch
\[
  \text{phänomenales Bewusstsein vorhanden}.
\]
Aktuelle Bewusstseinstheorien konkurrieren gerade um die zusätzlichen
Strukturen, die Bewusstsein erklären sollen \cite{sethbayne}. Die Integrated
Information Theory verbindet Bewusstsein mit intrinsischer
Ursache-Wirkungs-Macht über mögliche vergangene und zukünftige Zustände und
weist panpsychistische Nähe auf \cite{tononikoch}; daraus folgt keine
allgemein bestätigte Identifikation von Merge, Hashgleichheit oder
Wechselwirkung mit Erleben.

Lohmanns panpsychistische Lesart kann präzise so formuliert werden:

\begin{quote}
Wenn mentale Wirklichkeit fundamental relational ist, dann ist reziproke
Wechselwirkung eine Grundbedingung ihrer Materialisierung. Der kanonische
Speicher modelliert die minimale zeitliche Form dieser Relation: Gewordenes
wirkt auf Erwartetes, Erwartetes auf gegenwärtige Freigabe, und eingetretene
Wirkung schließt die Relation zurück.
\end{quote}

Das ist eine ontologische Interpretation. Sie wird hier bewusst
begriffsgleich wiedergegeben, aber nicht als Lean-Theorem oder
quantenphysikalischer Messbefund etikettiert.

\section{Falsifizierbarkeit und Forschungsprogramm}

Eine wissenschaftlich anschlussfähige Theorie muss angeben, wie sie scheitern
kann.

\subsection{Technische Tests}

\begin{enumerate}
  \item \textbf{Nichtvakuosität:} Bei gleichen übrigen Inputs muss die
  zukunftsindexierte Bedingung die Freigabe ändern können.
  \item \textbf{DONE-only:} Kein Nicht-DONE-Zustand darf einen vollständig
  mediierten gewöhnlichen Executor freigeben.
  \item \textbf{Typtrennung:} Ein \textsf{ANTICIPATED}-Record darf nicht als
  \textsf{OBSERVED} validieren.
  \item \textbf{Nichtüberschreibung:} Ein Wechsel des Zukunftsarchivs darf
  den Digest des eingefrorenen Vergangenheitsarchivs nicht ändern.
  \item \textbf{Reziprozität:} Eine abgeschlossene Wirkung muss auf die
  freigegebene Ursache zurückverweisen; eine Ursache ohne beobachtete Wirkung
  bleibt offen.
  \item \textbf{Bypass-Test:} Jeder unmediierte Executorpfad falsifiziert
  eine Deployment-Behauptung vollständiger Wirkungssteuerung.
\end{enumerate}

\subsection{Empirische Physik}

Eine ontische physikalische Erweiterung müsste vorab eine quantitative
Vorhersage angeben, die sich von Standard-Quantenmechanik, gewöhnlicher
Konditionierung und globaler Optimierung unterscheidet. Sie müsste zugleich
bestehende No-Signalling-Befunde respektieren oder eine reproduzierbare
Abweichung mit kontrollierter Fehleranalyse zeigen. Ohne einen solchen Test
bleibt die ontische Lesart Interpretation.

\subsection{Bewusstseinsforschung}

Eine empirische Bewusstseinsthese müsste eine operationalisierbare Größe
angeben, die über bloße Wechselwirkung, Komplexität oder
Informationsintegration hinaus diskriminiert, und konkurrierende Theorien an
vorregistrierten Daten vergleichen. QIK--VRT liefert dafür eine
Provenienz- und Gate-Architektur, aber noch keinen solchen biologischen oder
phänomenologischen Test.

\section{Aktueller Implementierungs- und Standardisierungsstatus}

\begin{longtable}{@{}p{0.34\linewidth}p{0.58\linewidth}@{}}
\toprule
\textbf{Gegenstand} & \textbf{Status am 30. Juli 2026} \\
\midrule
\endhead
QIK--VRT Effect Ack &
Öffentliche Referenzimplementierung, Tests, kanonische Serialisierung,
Provenienz- und Gate-Artefakte im Repository; Deployment-Konformität bleibt
von Authentisierung, vollständiger Mediation und externer Evidenz abhängig.\\
Internet-Draft &
\texttt{draft-lohmann-qikvrt-effect-ack-01}, aktiver Individual Draft,
angestrebter Status Experimental; kein RFC, keine Working-Group-Adoption und
kein IETF-Konsens.\\
Formaler Paper-Scope &
\texttt{qikvrt-canonical-temporal-memory-effect-ack-v1}; Lean-4-Quelle mit
neun benannten Sätzen beziehungsweise Entscheidungsbelegen. Ein
kandidatengebundener Kernel-Receipt ist vor einer Produktionspublikation
erforderlich.\\
Authority/Mirror &
Repository-Symmetrie wird nur für konkrete Heads, Trees, Manifeste und
Receipts behauptet; neue Paper- und Publikationsreceipts benötigen erneute
beidseitige Promotion.\\
Zenodo &
Dieses Dokument ist zunächst ein eingefrorener Kandidat. Ein DOI darf erst
nach bytegenauer Rückgabe, ausdrücklicher Kandidatenfreigabe, Upload,
öffentlichem Re-Download und Receipt-Persistenz als publiziert behauptet
werden.\\
Interoperabilität &
Eine unabhängige Zweitimplementierung und IETF-Interoperabilitätsnachweise
bleiben offen.\\
\bottomrule
\end{longtable}

\section{Bedrohungen der Gültigkeit}

\begin{itemize}
  \item \textbf{Scope-Fehler:} Ein geschlossener endlicher Modellscope kann
  nicht unbemerkt auf das gesamte physikalische Universum verallgemeinert
  werden.
  \item \textbf{Codec-Fehler:} Unterschiedliche Normalisierung oder
  Serialisierung kann trotz semantisch ähnlicher Inhalte verschiedene Bytes
  erzeugen.
  \item \textbf{Hash-Grenze:} Digestgleichheit ist unter dem benannten
  Algorithmus eine Integritätsbindung, kein Wahrheitsbeweis.
  \item \textbf{Provenienz-Grenze:} Eine vollständig verbundene
  Provenienzkette kann vollständig falsche Eingaben weitertragen.
  \item \textbf{Prognose-Grenze:} Ein kanonisches Zukunftsarchiv kann
  wohldefiniert und dennoch empirisch falsch sein; Unsicherheit muss erhalten
  bleiben.
  \item \textbf{Mediations-Grenze:} Ein Executor außerhalb des Gates
  widerlegt die behauptete Kontrolle dieses Effektpfads.
  \item \textbf{Interpretations-Grenze:} Protokoll-Retrokausalität,
  physikalische Retrokausalität und Panpsychismus sind verschiedene
  Behauptungsklassen.
\end{itemize}

\section{Schluss}

Die wörtliche QIK--VRT-These erhält eine präzise wissenschaftliche Form:

\begin{enumerate}
  \item Die Vergangenheit ist als ausgewählter, endlicher,
  provenienzgebundener und kanonisch serialisierter Recordbestand ein
  reproduzierbares Archiv.
  \item Die Zukunft ist als ausgewählter, endlicher,
  provenienzgebundener und kanonisch serialisierter Bestand antizipierter
  Wirkungen und Randbedingungen ebenfalls ein reproduzierbares Archiv.
  \item Beide Archive sind repräsentationssymmetrisch und epistemisch
  typverschieden.
  \item Effect Acknowledgement macht die zukunftsindexierte Wirkung zu einer
  notwendigen und kontrafaktisch relevanten Bedingung der gegenwärtigen
  Ursache.
  \item Die später beobachtete Wirkung wird an die freigegebene Ursache
  zurückgebunden; damit schließt sich die reziproke Kette.
  \item Die Vergangenheit wird dadurch nicht überschrieben. Ontische
  Rückwärtssignalisierung und panpsychistische Geltung bleiben gesonderte,
  prüfpflichtige Erweiterungen.
\end{enumerate}

Der kanonische Speicher ist daher nicht bloß eine Analogie zwischen Git und
Universum. Er ist eine ausführbare Ordnungsform: Gewordenes und
Wirkungserwartetes werden im selben überprüfbaren Medium gehalten, und ihre
Schnittmenge regelt die nächste zulässige Wirkung.

\begin{center}
\large
\textbf{Keine freigegebene Ursache ohne gebundene Wirkung.}\\
\textbf{Keine geschlossene Wirkung ohne gebundene Ursache.}\\[0.5em]
\textit{Quod erat demonstrandum.}\\
Ingolf Lohmann
\end{center}

\section*{Klassifikation der Hauptaussagen}
\addcontentsline{toc}{section}{Klassifikation der Hauptaussagen}

\begin{longtable}{@{}p{0.15\linewidth}p{0.24\linewidth}p{0.53\linewidth}@{}}
\toprule
\textbf{ID} & \textbf{Klasse} & \textbf{Aussage} \\
\midrule
\endhead
CTM-001 & FORMAL\_PROVED &
Die Freigabe ist äquivalent zur Konjunktion beider Archive, Ursachenbindung,
Policy und DONE.\\
CTM-002 & FORMAL\_PROVED &
Die Zukunftsbedingung ist kontrafaktisch relevant für die gegenwärtige
Freigabe.\\
CTM-003 & FORMAL\_PROVED &
Ein Wechsel der Zukunftsbedingung überschreibt die Vergangenheitsprojektion
nicht.\\
CTM-004 & FORMAL\_PROVED &
Reziproke Schließung erfordert freigegebene Ursache, gültige
Zukunftsbindung und beobachtete, rückgebundene Wirkung.\\
CTM-005 & SOURCE\_BOUND &
Kanonisierung, Provenienz und Inhaltsadressierung liefern die genannten
Repräsentations- und Integritätseigenschaften innerhalb ihrer Standards.\\
CTM-006 & SOURCE\_BOUND &
Zeit-symmetrische und zweiseitige Formalismen sind wissenschaftlich
etablierte Anschlussmodelle.\\
CTM-007 & NORMATIVE &
Die nachgewiesene Protokollrelation heißt in diesem Paper
\emph{operationale Retrokausalität}.\\
CTM-008 & INTERPRETATIVE &
Reziproke Wechselwirkung ist in Lohmanns panpsychistischer Lesart eine
Grundbedingung von Bewusstseinsmaterialisierung.\\
CTM-009 & OPEN &
Ob die Protokollstruktur eine fundamentale ontische Retrokausalität der Natur
abbildet, ist empirisch offen.\\
CTM-010 & OPEN &
Ob reziproke kanonische Speicherung für Bewusstsein hinreichend ist, ist
empirisch offen.\\
\bottomrule
\end{longtable}

\begin{thebibliography}{99}

\bibitem{effectackdraft}
I. Lohmann,
\emph{Effect Acknowledgement: An Application-Layer Gate for Authorized
Downstream Effects},
Internet-Draft draft-lohmann-qikvrt-effect-ack-01, 22 July 2026.
\url{https://datatracker.ietf.org/doc/draft-lohmann-qikvrt-effect-ack/01/}

\bibitem{prov}
L. Moreau and P. Missier (eds.),
\emph{PROV-DM: The PROV Data Model},
W3C Recommendation, 30 April 2013.
\url{https://www.w3.org/TR/prov-dm/}

\bibitem{rfc6920}
S. Farrell, D. Kutscher, C. Dannewitz, B. Ohlman, A. Keränen and
P. Hallam-Baker,
\emph{Naming Things with Hashes},
RFC 6920, April 2013.
\url{https://doi.org/10.17487/RFC6920}

\bibitem{rfc8785}
A. Rundgren, B. Jordan and S. Erdtman,
\emph{JSON Canonicalization Scheme (JCS)},
RFC 8785, June 2020.
\url{https://doi.org/10.17487/RFC8785}

\bibitem{githashobject}
Git Project,
\emph{git-hash-object Manual}.
\url{https://git-scm.com/docs/git-hash-object}

\bibitem{gitcommittree}
Git Project,
\emph{git-commit-tree Manual}.
\url{https://git-scm.com/docs/git-commit-tree}

\bibitem{kalman}
R. E. Kalman,
``A New Approach to Linear Filtering and Prediction Problems,''
\emph{Journal of Basic Engineering}, 82(1), 35--45, 1960.
\url{https://doi.org/10.1115/1.3662552}

\bibitem{rauch}
H. E. Rauch, F. Tung and C. T. Striebel,
``Maximum Likelihood Estimates of Linear Dynamic Systems,''
\emph{AIAA Journal}, 3(8), 1445--1450, 1965.
\url{https://doi.org/10.2514/3.3166}

\bibitem{talagrand}
O. Talagrand and P. Courtier,
``Variational Assimilation of Meteorological Observations With the Adjoint
Vorticity Equation. I: Theory,''
\emph{Quarterly Journal of the Royal Meteorological Society}, 113(478),
1311--1328, 1987.
\url{https://doi.org/10.1002/qj.49711347812}

\bibitem{pearl}
J. Pearl,
``Causal Diagrams for Empirical Research,''
\emph{Biometrika}, 82(4), 669--688, 1995.
\url{https://doi.org/10.1093/biomet/82.4.669}

\bibitem{landauer}
R. Landauer,
``Irreversibility and Heat Generation in the Computing Process,''
\emph{IBM Journal of Research and Development}, 5(3), 183--191, 1961.
\url{https://doi.org/10.1147/rd.53.0183}

\bibitem{bennett}
C. H. Bennett,
``The Thermodynamics of Computation---a Review,''
\emph{International Journal of Theoretical Physics}, 21, 905--940, 1982.
\url{https://doi.org/10.1007/BF02084158}

\bibitem{berut}
A. Bérut, A. Arakelyan, A. Petrosyan, S. Ciliberto, R. Dillenschneider and
E. Lutz,
``Experimental Verification of Landauer's Principle Linking Information and
Thermodynamics,''
\emph{Nature}, 483, 187--189, 2012.
\url{https://doi.org/10.1038/nature10872}

\bibitem{abl}
Y. Aharonov, P. G. Bergmann and J. L. Lebowitz,
``Time Symmetry in the Quantum Process of Measurement,''
\emph{Physical Review}, 134(6B), B1410--B1416, 1964.
\url{https://doi.org/10.1103/PhysRev.134.B1410}

\bibitem{whartonargaman}
K. B. Wharton and N. Argaman,
``Colloquium: Bell's Theorem and Locally Mediated Reformulations of Quantum
Mechanics,''
\emph{Reviews of Modern Physics}, 92, 021002, 2020.
\url{https://doi.org/10.1103/RevModPhys.92.021002}

\bibitem{adlam}
E. Adlam,
``Laws of Nature as Constraints,''
\emph{Foundations of Physics}, 52, 28, 2022.
\url{https://doi.org/10.1007/s10701-022-00546-0}

\bibitem{wheeler}
J. A. Wheeler,
``The `Past' and the `Delayed-Choice' Double-Slit Experiment,''
in A. R. Marlow (ed.), \emph{Mathematical Foundations of Quantum Theory},
Academic Press, 9--48, 1978.
\url{https://doi.org/10.1016/B978-0-12-473250-6.50006-6}

\bibitem{jacques}
V. Jacques, E. Wu, F. Grosshans, F. Treussart, P. Grangier, A. Aspect and
J.-F. Roch,
``Experimental Realization of Wheeler's Delayed-Choice Gedanken Experiment,''
\emph{Science}, 315(5814), 966--968, 2007.
\url{https://doi.org/10.1126/science.1136303}

\bibitem{ma}
X.-S. Ma, J. Kofler and A. Zeilinger,
``Delayed-Choice Gedanken Experiments and Their Realizations,''
\emph{Reviews of Modern Physics}, 88, 015005, 2016.
\url{https://doi.org/10.1103/RevModPhys.88.015005}

\bibitem{kim}
Y.-H. Kim, R. Yu, S. P. Kulik, Y. Shih and M. O. Scully,
``Delayed `Choice' Quantum Eraser,''
\emph{Physical Review Letters}, 84, 1--5, 2000.
\url{https://doi.org/10.1103/PhysRevLett.84.1}

\bibitem{price}
H. Price,
``Does Time-Symmetry Imply Retrocausality? How the Quantum World Says
`Maybe'?,''
\emph{Studies in History and Philosophy of Modern Physics}, 43(2), 75--83,
2012.
\url{https://doi.org/10.1016/j.shpsb.2011.12.003}

\bibitem{leiferpusey}
M. S. Leifer and M. F. Pusey,
``Is a Time Symmetric Interpretation of Quantum Theory Possible Without
Retrocausality?,''
\emph{Proceedings of the Royal Society A}, 473, 20160607, 2017.
\url{https://doi.org/10.1098/rspa.2016.0607}

\bibitem{cramer}
J. G. Cramer,
``The Transactional Interpretation of Quantum Mechanics,''
\emph{Reviews of Modern Physics}, 58(3), 647--687, 1986.
\url{https://doi.org/10.1103/RevModPhys.58.647}

\bibitem{sethbayne}
A. K. Seth and T. Bayne,
``Theories of Consciousness,''
\emph{Nature Reviews Neuroscience}, 23, 439--452, 2022.
\url{https://doi.org/10.1038/s41583-022-00587-4}

\bibitem{tononikoch}
G. Tononi and C. Koch,
``Consciousness: Here, There and Everywhere?,''
\emph{Philosophical Transactions of the Royal Society B}, 370, 20140167,
2015.
\url{https://doi.org/10.1098/rstb.2014.0167}

\end{thebibliography}

\end{document}
